% Chapter 8: Conclusion
\section{Conclusão e Trabalhos Futuros}

Este projeto de TCC I delineou a concepção e o planejamento do jogo sério "Alchemon", uma ferramenta destinada a apoiar o ensino da Tabela Periódica. Foi realizada a fundamentação teórica, a definição do problema, dos objetivos e da metodologia de desenvolvimento. O design do jogo foi estruturado para integrar de forma coesa os conteúdos pedagógicos com mecânicas de jogo engajadoras, seguindo as melhores práticas da área de aprendizagem baseada em jogos \cite{book_example, journal_example}.

O trabalho realizado até aqui corresponde às fases iniciais do ciclo de desenvolvimento, resultando em um plano sólido para a implementação do protótipo. Acreditamos que a abordagem proposta tem grande potencial para criar uma experiência de aprendizagem positiva e eficaz.

\subsection{Trabalhos Futuros}
A continuação deste projeto, a ser realizada no TCC II, envolverá as seguintes etapas:
\begin{itemize}
    \item Conclusão do desenvolvimento do protótipo funcional do "Alchemon".
    \item Realização de testes de usabilidade para refinar a interface e a experiência do usuário.
    \item Implementação de um estudo de caso em um ambiente educacional real, com a participação de estudantes do ensino médio.
    \item Coleta e análise de dados quantitativos e qualitativos para avaliar o impacto do jogo na motivação e no desempenho acadêmico dos alunos.
    \item Documentação dos resultados e redação das conclusões finais da pesquisa.
\end{itemize}

Espera-se que, ao final do projeto completo, o "Alchemon" não seja apenas um produto final, mas também uma contribuição validada para o campo da tecnologia educacional, como sugerido em conferências sobre inovação \cite{conference_example}.
